{\bfseries RSA: cifrado y descifrado. Se eligen dos primos pequeños $p = 7$ y $q = 13$.}

\begin{enumerate}[a)]
	\item {\bfseries Calcule $n = p \cdot q$ y $\varphi(n) = (p-1)(q-1)$.}

	\vspace{3cm}

	\item {\bfseries Elija un exponente público $e$ tal que $1 < e < \varphi(n)$ y $\gcd(e, \varphi(n)) = 1$. Tome $e = 5$. Justifique que es válido.}

	\vspace{3cm}

	\item {\bfseries Calcule el exponente privado $d = e^{-1} \bmod \varphi(n)$ usando el algoritmo extendido de Euclides (muestre los pasos).}

	\vspace{5cm}

	\item {\bfseries El mensaje a cifrar es $M = 20$. Calcule el criptograma $C = M^e \bmod n$ (use exponenciación modular rápida, mostrando las potencias necesarias).}

	\vspace{5cm}

	\item {\bfseries Descifre $C$ calculando $C^d \bmod n$ y verifique que obtiene $M = 20$.}

	\vspace{5cm}

	\item {\bfseries ¿Cuál es la clave pública y cuál la privada?}

	\vspace{3cm}

\end{enumerate}
